%% ma: L10-B13-0002
%% lop: 10
%% chuong: 5
%% bai: 13
%% dang: 10.13.2
%% loai: TL
%% muc: TH
%% dap_an: "-"
%% nguon: "chuyen-de-cac-so-dac-trung-cua-mau-so-lieu-khong-ghep-nhom-toan-10-chuong-trinh-moi.pdf (GV. Phan Nhật Linh) (P02), Bài 02 Dạng 1 Bài tập tự luận 2"
%% kien_thuc: "Lớp 10 Bài 13 (số trung bình)"
%% trang_thai: chinh-thuc
%% ghi_chu: "Tự duyệt theo yêu cầu giáo viên 22/10/2026."
\begin{ex}
Bảng sau cho biết thời gian chạy cự li $100$ m của các bạn trong lớp (đơn vị: giây):
\begin{center}
\begin{tabular}{|c|c|c|c|c|c|}\hline
Thời gian & $12$ & $13$ & $14$ & $15$ & $16$\\\hline
Số bạn & $5$ & $7$ & $10$ & $8$ & $6$\\\hline
\end{tabular}
\end{center}
Hãy tính thời gian chạy trung bình cự li $100$ m của các bạn trong lớp (làm tròn đến hàng phần nghìn).
\loigiai{$n=5+7+10+8+6=36$.

$\overline x=\dfrac{5\cdot12+7\cdot13+10\cdot14+8\cdot15+6\cdot16}{36}=\dfrac{507}{36}\approx14{,}083$ (giây).}
\end{ex}
